\chapter{Society and Governance Parallels in Minecraft: The Systems Experiment Findings}

\begin{center}
    Author: Mila. Editors: Harlow and retrooper. Peer reviewed. % clearly indicate non-peer reviewed articles!
\end{center}

\begin{flushright}
    IMSC: j(2)-mia-1

    DOI: \doilink{10.5281/zenodo.23702314}
\end{flushright}

% ============================================================================
\section*{Abstract}
The Minecraft world is an example of a digital space that closely resembles our real world. TSE specifically, just like many other multiplayer servers, is an ideal place to observe said similarities. By finding connections between both worlds involving phenomena such as how power is distributed, diplomacy handled, or how the world progresses, we can also use this information to make predictions about events involving those topics. This can be achieved through a simple process involving a little bit of history and logic.

\section{The commons of different worlds}
In a world where the internet bridges large distances and connects individuals throughout the globe, there are bound to be certain applications that will simulate the world itself. Games like Minecraft that allow players to gather and build cultures and governments represent this perfectly; however, many people fail to realize how similar a simple game can match our large and complex society. In the instance of the Systems Experiment server, a multiplayer Minecraft server with two opposing “countries”, three particular topics stood out to strongly resemble their real-world counterparts. By gathering the perspectives of the majority of individuals per team and contrasting them with movements and popular opinions in the real-world, the resemblance between a Minecraft society and our real-world societies is evident. By the end of the experiment, many players witnessed the same entropy that came with technology, as well as many diplomatic events that could only be synthesized by one’s experiences in the real-world. With these parallels drawn, people can attempt to explain real-world phenomena, explore the causes of important events, and even make their own predictions about the world we live in. Although power is a highly debated topic, for the sake of this article, we can describe power as how capable a person is of changing the behaviors of a group of people. In our world, power is distributed and held in many different ways, with examples ranging from influencers on social media, to people with enough power to conceal themselves from public knowledge. One concentration of power, however, is visible in both TSE and real life—that being politicians. In real life, politicians obtain power through public support, control of resources, and by utilizing their ability to change government policy. To determine whether the leadership roles in the IMS experiment match those of politicians, we must draw connections between these aspects.

To start, within the experiment server, both teams required their leaders to hold a sufficient amount of support from their people. This can be particularly seen with Team Alpha’s first election on July 11, 2026, where Phantom won over Scorvon in a poll. This was achieved partially in his speech, where his promises for stability and his anti-war emphasis reflected the concerns of the people, similar to real life campaigns. Actions also are involved, as some voters might have cast doubt on Scorvon’s own pro-diplomatic claims after overhearing his threats toward Ordonia. Next, itistiff and Llama Amongus, in charge of Team Ordonia’s relations with Team Alpha, were able to permit the allocation of resources to provide trade to team Alpha. We will soon see why they did this, but said trade still holds a considerable amount of diplomatic power, as its use (or lack of it) can determine friends, foes, or everything in between. In instances of changing government policy, leaders like Scorvon, Phantom, and Pheelix signed agreements espionage, trade, and etiquette. At its core, agreements like these limit or expand the freedoms of players/people for a key change – mutual benefit. This benefit can boost both teams in many ways, while strengthening the relationship between the leaders and the populations, which is also useful for diplomacy. With these connections made, we can safely conclude that there is at least one major similarity in powerholders between the real world and TSE, however, we can also speculate about several other sources of power between the two worlds.

An apparent pattern that occurs in almost any world is material wealth. Especially in games like Minecraft where “grinding” can be used to obtain resources, people often resort to obtaining a surplus of materials or currency to gain power. With the use of material, this power can come in different forms, whether that be militaristic power (hiring soldiers/buying gear) or political power (lobbying/handouts). Overall, since power (or its use of it) can be a major factor in deciding actions or consequences in both worlds, this will aid the connection between the two worlds as well as predictions regarding politics.

Although wars in Minecraft are by no means the same as those in real life, the similar effects of war (death, destruction, loss of materials) makes diplomacy strikingly similar in both worlds. For a base definition of diplomacy, we could say that countries use communication, compromise, and peacebuilding to prevent conflicts and resolve incidents. Although partially kept secret in TSE, these same three tactics could be seen being used by the leaders of team Alpha, the leader of team Ordonia, and the diplomatic board of team Ordonia. To start, by interviewing a couple of members of Ordonia (who I will keep anonymous in this article), they have admitted to using trade deals to try and prevent Alpha from wanting a war. This is a good example of peacebuilding, and a popular real-world tactic known as appeasement. Similarly to countries utilizing appeasement, Ordonia simply did not have the manpower or the resources to compete with Alpha militaristically, leading them to resort to trade to prevent the war entirely. There were also instances of meetings where both teams communicated about trades, issues, and the war they were trying to prevent. These were rather popular, to the extent that Team Alpha built a meeting place in their town hall to host them. Notably, it was both large and decorated, either to impress Ordonian diplomats and any spectators, or to ensure safety and convenience when hosting a meeting. This also took place to a lesser extent in smaller visits like tea parties, which were not exactly diplomatic oriented, but still served as proper communication and as an opportunity to build relations between teams at a personal level, which cannot be easily achieved with deals or diplomatic visits.

There have also been instances of compromise taking place within these meetings, specifically about policies and agreements that the two teams have created and upheld throughout the server. A good example of this would be how team Ordonia compromised to let team Alpha enter their territory, on the condition that any Alpha players visiting would announce themselves and keep their locator bar on. In terms of this agreement, it was much less implemented as a way to avoid war than a simple courtesy for Team Alpha as a whole. Courtesies such as this compromise can greatly affect opinions and relations, so they occasionally will be performed to improve one’s image instead of genuine goodwill. As the event starts to come to a close (as of writing this article), it appears that these methods of diplomacy have worked, as a war has not started, and there is not any motive or motivation to start one in the remaining time. This connection can also serve a good purpose, as being able to predict wars and fights will aid the understanding of the players and the server better.

As mentioned in the introduction of this article, there has been an observable entropy linked with in-game progression, which is similar to the side effects we see with technology advancement in the real-world. To understand this fully, we must first look at how our world was affected by technology. At the surface level, it has improved our lives – people can live comfortably in soft mattresses and fancy smart homes, with technological efficiency in almost every corner of the world. The hard, dirty jobs have been replaced by large machinery that can automate grueling tasks at the press of a button, and discoveries like flight and electricity have made accomplishments that our ancestors could not even dream of. This is comparable to advancements in Minecraft too, with the unlocking of new resources like netherite, improvements in farms such as trading halls, and new items such as elytras. The most obvious drawback to new technology and improvements is the improvement of war.

Wars can be fought with brutal efficiency, with machine guns consecutively killing hundreds of people, and bombs being dropped to raze cities in mere seconds. Again, the same exists in Minecraft. Maces and wind charges can deal damage that no other weapon can, and massive redstone machines can be sent through the map, blowing up everything in their way. To most, this would seem as an appropriate consequence for the progression of technology, however, there are still much darker consequences that little realize until it is too late.

\section{Finding use in patterns}
With the evolution of society, depression, and other mental health issues have been rapidly increasing.~\citep{Hidaka(2012)} People’s lives have been watered down to 9-5 jobs, with the rest of people’s time spent idly watching TV, playing games, or being on social media.~\citep{Malinsky(2026)} This too is present in Minecraft, particularly when people get bored and burn out in the later stages of a server, rather than in the start of a fresh server. After all, this mainly happens after there is nothing new to do, or when people lose motivation to continue playing. A part of this is that the same technology -- that sought to improve lives and automate the game -- strips the joy and risk out of simple adventures and daily activities. In this regard, machines like auto miners or raid farms automatically harvest resources rather than having the player earn them, making it a repetitive clicking session instead of a thrilling outing. People can also burn out from the work, as the bigger end-game projects can be the same repetitive grinding we see in real life jobs. Several players fell victim to this in TSE. For example, one player (who will remain anonymous) experienced burnout from trying to figure out late-game redstone. But at the end of the day, sometimes there are still simple pleasures. After all, there is nothing quite like having a simple cup of tea in Minecraft, or in real-life.

Now that several key connections have been established, they can be used to predict certain related aspects about both Minecraft and real-life. To perform this, simply start by viewing a scenario in a “cause and effect” sort of way. Although being a simplified process, if supplemented with the correct data, it can utilize the past to provide a result. The past, being a useful tool, can be used to see many human behavioral patterns. Human history particularly is also one of the largest and most accessible pasts, making it accessible and plentiful. For a hypothetical scenario, let's start with:

\begin{verse}
    Cause: Tensions are very high between team Alpha and team Ordonia Effect:

    ???
\end{verse}

To fill in the effect, all we would need to do is ask ourselves “if two countries had high tensions, what would they do in real life?” Then, using our historical knowledge, or knowledge about the world in the present world, we can fill in:

\begin{verse}
    Cause: Tensions are very high between team Alpha and team Ordonia Effect:

    Both teams will attempt to use diplomacy to ease the tensions.
\end{verse}

Obviously, we cannot be certain if this is accurate at face value. If the cause is more specific (e.g., motives, amounts of resources already, etc.), the effect will also be more specific, and with predictions, even if you had a surplus of data for your cause, you can never be 100\% sure about your effect.

\section{Conclusion}
All in all, the connections we can draw between the world of TSE and the physical world and the subsequent predictions only mark the start of what we can achieve with pattern recognition and an observant eye. Looking closely, it’s also not hard to see common themes like power, diplomacy, and technology in other player-driven games too. However, just remember that even the comforting, blocky landscape of Minecraft can still be strikingly similar to the outside world we all know.

\section{Declarations}
The author and the reviewer have declared a conflict of interest, having participated within The Systems Experiment. The editor Harlow has declared a conflict of interest, having played and organized The Systems Experiment. No conflict of interest has been declared by the editor retrooper. No AI use has been declared.

    {
        \raggedright
        \begin{thebibliography}{99}
            \bibitem[Malinsky(2026)]{Malinsky(2026)}
            Malinsky, Gili. “The No. 1 Way Americans Spend Their Free Time: It's ‘Always Been a Huge Part’ of Our Culture.” CNBC, 23 July 2026, https://www.cnbc.com/2026/07/23/heres-the-no-1-way-americans-spend-their-leisure-time.html. Accessed 27 July 2026.
            \bibitem[Hidaka(2012)]{Hidaka(2012)}
            Hidaka, Brandon H. “Depression as a Disease of Modernity: Explanations for Increasing Prevalence.” Journal of Affective Disorders, vol. 140, no. 3, Nov. 2012, pp. 205–214, https://pmc.ncbi.nlm.nih.gov/articles/PMC3330161/, 10.1016/j.jad.2011.12.036. Accessed 6 August 2026.

        \end{thebibliography}
    }